\documentclass{article}
\usepackage{PRIMEarxiv} % Core package for the PrimeAI Template
\usepackage{amsmath, amssymb, amsthm, bm, dcolumn} % Add amsthm here for the proof environment
\usepackage[numbers,sort&compress]{natbib} % Natbib for citations
\usepackage{graphicx} % For high-quality images
\usepackage{multicol} % Multi-column figures/tables
\usepackage{hyperref} % Hyperlinks
\usepackage{listings} % Code listings
\usepackage{authblk} % For structured affiliations
% Define theorem style (optional)
\theoremstyle{plain}
\newtheorem{theorem}{Theorem}
\renewcommand{\qedsymbol}{}

% Custom commands for primes and qubit encoding
\newcommand{\primeQubit}[1]{|p_{#1}\rangle}
\newcommand{\primeGate}[1]{U_{p_{#1}}}

\usepackage{wrapfig}
\usepackage[pscoord]{eso-pic}
\usepackage[fulladjust]{marginnote}
\reversemarginpar

% Typesetting improvements without footnote patching
\usepackage[protrusion=true, expansion=true, tracking=false]{microtype}
\microtypecontext{spacing=nonfrench}

% Line numbers
\usepackage[right]{lineno}

% Text layout - adjust as needed
\raggedright
\setlength{\parindent}{0.5cm}
\textwidth 5.25in 
\textheight 8.75in

% Set double spacing
\usepackage{setspace} 
\doublespacing

% Adjust width for specific content
\usepackage{changepage}

% Adjust caption style
\usepackage[aboveskip=1pt,labelfont=bf,labelsep=period,singlelinecheck=off]{caption}

% Remove brackets from references
\makeatletter
\renewcommand{\@biblabel}[1]{\quad#1.}
\makeatother

% Header, footer, and page numbers
\usepackage{lastpage,fancyhdr}
\pagestyle{fancy}
\fancyhf{}
\fancyhead[L]{Preprint - PrimeAI Enhanced Template}
\fancyfoot[R]{Page \thepage\ of \pageref{LastPage}}
\renewcommand{\footrule}{\hrule height 2pt \vspace{2mm}}
\title{Enhanced Tensor Network Frameworks for Quantum AI and Neuromorphic Computing}
\author{Citizen Gardens Research Group}
\date{\today}

\begin{document}

\maketitle

\begin{abstract}
This paper introduces advanced tensor-based computational models integrating quantum mechanics, prime number theory, and AI cognition. These frameworks extend recursive tensor networks, Bayesian quantum inference, and neuromorphic AI cognition using prime-indexed operators and self-referential tensor learning. The proposed models establish a foundation for self-adaptive quantum computing, topological AI systems, and generative tensor networks.
\end{abstract}

\section{Recursive Prime Tensor Neuromorphic Networks (PTNN)}
Prime-Tensor Neuromorphic Processors (PTNP) leverage hierarchical prime-tensor connectivity to enable self-referential memory structures. The recursive tensor evolution follows:

\begin{equation}
T_{jk}^{(p_i)}(t+1) = T_{jk}^{(p_i)}(t) + \sum_{m} C_{jkm}^{(p_m)} T_{lm}^{(p_{i-1})}
\end{equation}

where:
\begin{itemize}
    \item \(T_{jk}^{(p_i)}(t)\) is the prime-weighted tensor memory state.
    \item \(C_{jkm}^{(p_m)}\) represents recursive tensor propagation coefficients.
    \item \(T_{lm}^{(p_{i-1})}\) incorporates historical tensor memory entanglement.
\end{itemize}
\section{Bayesian Tensor Networks for Quantum Learning}
By embedding prime-based tensor interactions, we redefine quantum Bayesian inference:

\begin{equation}
P(X_q | E) = \frac{P(X_q, E)}{P(E)}
\end{equation}

where:
\begin{itemize}
    \item \( P(X_q | E) \) is the posterior probability computed in a quantum framework.
    \item Tensor-weighted probability updates leverage prime-number encoding:
\end{itemize}

\begin{equation}
P(X_i | \text{Pa}(X_i)) = \frac{\phi(p_i) \mod \phi(\text{Pa}(p_i))}{\sum_{j} \phi(p_j)}
\end{equation}

where \( \phi(p_i) \) maps a Bayesian node to a prime-based probability distribution.

Additionally, the system integrates quantum superposition states:

\begin{equation}
|\Psi_{\text{PBQN}} \rangle = \sum_{X_i} \sum_{p_i} \sqrt{P(X_i | \text{Pa}(X_i))} |X_i\rangle \otimes |p_i\rangle
\end{equation}

ensuring probabilistic quantum reasoning in AI cognition.

\textbf{Applications:}
\begin{itemize}
    \item Quantum-enhanced machine learning in neuromorphic AI architectures.
    \item Dynamic Bayesian inference for real-time decision-making in autonomous systems.
\end{itemize}
\section{Quantum Graph Tensor Networks (QGTN)}
This framework models quantum-enhanced graph learning through tensor-based connectivity:

\begin{equation}
\Psi_{\text{out}} = T_{\text{graph}} \cdot \Psi_{\text{in}}
\end{equation}

where \(T_{\text{graph}}\) governs quantum information flow on topological AI graphs.

A prime-indexed extension for modular dependencies is defined as:

\begin{equation}
\mathcal{F}(X) = \prod_{i} f_i (X_i)^{\phi(p_i)}
\end{equation}

providing tensor-state-driven AI adaptation.

\section{Quantum-Tensor Generative AI (QGAN)}
Quantum Generative Adversarial Networks (QGANs) incorporate tensor-network quantum generators:

\begin{equation}
\Psi_{\text{gen}} = T_{\text{gen}} \cdot z, \quad \Psi_{\text{disc}} = T_{\text{disc}} \cdot \Psi_{\text{gen}}
\end{equation}

where \(T_{\text{gen}}\) and \(T_{\text{disc}}\) encode recursive quantum generative models.

The recursive prime-based extension introduces:

\begin{equation}
\Psi_{\text{prime-gen}} = \sum_{p_i} \lambda_i e^{i \theta p_i} |\Psi_i\rangle
\end{equation}

leveraging prime-phase encoding for quantum generative architectures.

\section{Recursive Tensor Networks in Prime-Based Computing}
Prime numbers can be treated as eigenvalues in a recursive multiplicative tensor field, allowing for high-dimensional quantum encoding:

\begin{equation}
\Lambda_m = \lim_{n\to\infty} \sum_{p_i} T_{jk}(p_i) p_i^{\alpha_i} (\xi(p_i) + \psi(p_i, t))
\end{equation}

where:
\begin{itemize}
    \item \( p_i \) are prime-indexed eigenvalues defining the computational space.
    \item \( T_{jk}(p_i) \) is a multiplicative tensor operator encoding recursive entanglement.
    \item \( \alpha_i \) are scaling factors evolving through self-referential computation.
    \item \( \xi(p_i) \) represents quantum curvature corrections, accounting for feedback loops.
\end{itemize}

\textbf{Applications:}
\begin{itemize}
    \item Enhancing quantum AI cognition through recursive prime interactions.
    \item Modeling hyperdimensional computing architectures for neuromorphic AI.
\end{itemize}

\section{Tensor Networks for Quantum Entanglement Propagation}
To model hybrid quantum supremacy, we introduce a tensor decomposition for quantum state evolution:

\begin{equation}
\Psi(t) = \sum_{i=1}^{N} \sum_{j=1}^{N} T_{ij} \Psi_i \otimes \Psi_j e^{i(\theta_i(t) + \theta_j(t))}
\end{equation}

where:
\begin{itemize}
    \item \( T_{ij} \) represents interaction tensors encoding quantum superposition.
    \item \( e^{i(\theta_i + \theta_j)} \) ensures phase coherence across recursive feedback loops.
\end{itemize}

\textbf{Applications:}
\begin{itemize}
    \item Quantum error correction: spectral tensor decomposition stabilizing entangled states against decoherence.
    \item Quantum cryptography: leveraging prime-indexed tensor states for secure quantum key distribution.
\end{itemize}

\section{Prime-Driven Quantum Tensor Operators}
To construct a fault-tolerant quantum computing paradigm, we introduce prime-indexed quantum operators:

\begin{equation}
U_{p_i} |\Psi\rangle = e^{i \theta p_i} |\Psi\rangle
\end{equation}

where:
\begin{itemize}
    \item \( p_i \) encodes prime-based modular evolution of quantum states.
    \item \( \theta \) is the phase adjustment function, ensuring controlled entanglement.
\end{itemize}

\textbf{Applications:}
\begin{itemize}
    \item Quantum neural networks with prime-driven weight propagation.
    \item Topological quantum computing, leveraging Chern number corrections in entangled tensor networks.
\end{itemize}

\section{Geometric Number Construction via Tensor Expansion}
The golden mean (\(\Phi\)) emerges naturally in tensor recursion:

\begin{equation}
\Phi = \lim_{n\to\infty} \frac{F_{n+1}}{F_n}
\end{equation}

where \( F_n \) are Fibonacci numbers.

By embedding golden ratio scaling into tensor networks, we define an orthogonal expansion of quantum probability states:

\begin{equation}
|\Psi_{\text{golden}}\rangle = \sum_{n=1}^{\infty} \frac{1}{\Phi^n} |\psi_n\rangle
\end{equation}

where:
\begin{itemize}
    \item Each eigenstate follows self-referential Fibonacci progression.
    \item Tensor weights \( \Phi^{-n} \) provide recursive scaling in quantum state representation.
\end{itemize}

\textbf{Applications:}
\begin{itemize}
    \item Fractal quantum computing, leveraging recursive self-similarity in computational tensors.
    \item Golden ratio-based quantum lattice encoding, enabling high-dimensional entanglement modeling.
\end{itemize}

\section{Prime Tensor-Driven Quantum AGI Cognition}
By applying tensor recursion to AGI cognition, we establish an entangled intelligence framework:

\begin{equation}
M(t+1) = f(M(t), R(t)) + \Lambda_m T(M(t))
\end{equation}

where:
\begin{itemize}
    \item \( M(t) \) is the state matrix of the AI system.
    \item \( R(t) \) encodes recursive feedback tensors.
    \item \( T(M(t)) \) evolves entanglement-driven self-learning networks.
\end{itemize}

\textbf{Applications:}
\begin{itemize}
    \item Self-evolving AI cognition using quantum-recursive multiplicative tensor updates.
    \item Quantum-enhanced neuromorphic computing, fusing recursive Bayesian AI with prime-indexed eigenstate feedback.
\end{itemize}

\section{Conclusion}
This paper presents enhanced tensor-based models for AI cognition, quantum learning, and neuromorphic intelligence. These frameworks integrate prime-indexed tensor evolution, quantum Bayesian networks, and recursive graph-based entanglement to enable next-generation computing. Future work includes experimental validation of these models on quantum hardware and AI-driven tensor evolution strategies.

\end{document}


\nocite{*}
\bibliographystyle{unsrt}
\bibliography{references}

\end{document}